% generated by roam.py handbook export -- не редактировать, править заметки
% 20261010152131-сходимость_градиентного_спуска_с_шагом_армихо_вульфа.org (body 9ad7f26f3d96)

\begin{theorem}[{Сходимость градиентного спуска с шагом Армихо-Вульфа}]\label[theorem]{note:36F6E1BA-41DB-460D-A9C6-B980621727D4}
  Пусть:

  \begin{itemize}
    \item $f \in C_L^{1,1}$ (\hyperref[note:7DD1C7AB-1A49-4675-9F21-5A16C717F670]{Функции с липшицевым градиентом})
    \item $x_{k+1} = x_k + \alpha_k d_k$, где $\left\langle \nabla f(x_k); d_k \right\rangle < 0$~--- \hyperref[note:DF0CC2BA-9234-47D3-8988-7A1DAFC9B114]{градиентный спуск с шагом Армихо-Вульфа}
    \item $\exists f^{*}$~--- минимальное значение функции
    \item Зафиксируем $0<c_{1}<c_2<1$ из \hyperref[note:5330B23D-95DD-4D5D-B4A3-63EBADB1DB7D]{Условия Армихо-Вульфа}
    \item $\exists C >0:\ |\cos \theta_k| > C$ (из \hyperref[note:93FC9CC4-79FA-4CEA-9B10-9DBFB2C07850]{Оценка на шаг Армихо-Вульфа для градиентного спуска})
    \item $\alpha_k$~--- подчиняется \hyperref[note:5330B23D-95DD-4D5D-B4A3-63EBADB1DB7D]{Условия Армихо-Вульфа}
  \end{itemize}

  Тогда $\|\nabla f(x_k)\|^2 \rightarrow 0$ и верна такая оценка:
  \[
  \min_{0\le k \le n} \| \nabla f(x_k)\| \le \sqrt{\frac{L (f(x_0) - f^{*})}{(n+1)c_1(1-c_2)C^2}}
  \]
\end{theorem}

\begin{remark}
  При оптимальных $c_1,c_2,C$ оценка будет стремиться к
  \[
  \min_{0\le k \le n} \|\nabla f(x_k)\| \le \sqrt{\frac{4L (f(x_0) - f^{*})}{n+1}},
  \]
  что в $\sqrt{2}$ раза хуже, чем у \hyperref[note:C3B64D64-D35B-42AC-B683-6736227CBBBC]{спуска с постоянным шагом} $\frac{1}{L}$. Зато в нашем случае не обязательно знать $L$.
\end{remark}

\begin{proof}
  Хотим доказать, что
  \[
  f(x_{k+1}) \le f(x_k) - \mathrm{const} \cdot \|\nabla f(x_k)\|^2
  \]
  А дальше с помощью трюка с телескопической суммой из \hyperref[note:C3B64D64-D35B-42AC-B683-6736227CBBBC]{Сходимость градиентного спуска для функций с липшицевым градиентом} можно вывести такую оценку;

  Заметим, что \hyperref[note:15611E05-13FC-4FB0-AE00-A4BB857A29E6]{Условие Армихо} влечет за собой оценку (Это описано в \hyperref[note:DF0CC2BA-9234-47D3-8988-7A1DAFC9B114]{Градиентный спуск с шагом Армихо-Вульфа})
  \[
  f(x_{k+1}) \le f(x_k) + c_1\alpha_k \left\langle \nabla f(x_k);d_k \right\rangle
  \]

  Давайте выпишем \hyperref[note:93FC9CC4-79FA-4CEA-9B10-9DBFB2C07850]{Оценка на шаг Армихо-Вульфа для градиентного спуска} для $\alpha_k$
  \[
  \alpha_k \ge \frac{(c_2-1)\|\nabla f(x_k)\|}{L \|d_k\|} \cdot \cos \theta_k
  \]
  Так как $\left\langle \nabla f(x_k);d_k \right\rangle < 0$ (это обсуждалось в \hyperref[note:DF0CC2BA-9234-47D3-8988-7A1DAFC9B114]{Градиентный спуск с шагом Армихо-Вульфа}), то
  уменьшение параметра $\alpha$ перед ним увеличит значение (т.к. число отрицательное).
  тогда
  \begin{multline*}
  f(x_{k+1}) \le f(x_k) + c_1\alpha_k \left\langle \nabla f(x_k);d_k \right\rangle \le \\
  f(x_k) + \frac{c_1(c_2-1)\|\nabla f(x_k)\|}{L \|d_k\|} \cdot \cos \theta_k \cdot \left\langle \nabla f(x_k); d_k \right\rangle
  \end{multline*}
  Заметим, что
  \[
  \frac{1}{\|d_k\|} \cdot \left\langle \nabla f(x_k); d_k \right\rangle = \cos \theta_k \cdot \|\nabla f(x_k)\|
  \]
  Тогда
  \begin{multline*}
  f(x_{k+1}) \le
  f(x_k) + \frac{c_1(c_2-1)\|\nabla f(x_k)\|}{L \|d_k\|} \cdot \cos \theta_k \cdot \left\langle \nabla f(x_k); d_k \right\rangle = \\
  f(x_k) + \frac{c_1(c_2-1)\cos^2 \theta_k \|\nabla f(x_k)\|^2}{L}
  \end{multline*}
  Воспользуемся ограниченностью $\cos \theta_k$ снизу. Оценка верхняя так как $c_2 - 1 < 0$, следовательно вычитать мы будем меньше. Получим
  \[
  f(x_{k+1}) \le f(x_k) + \frac{c_1(c_2-1)C^2}{L} \cdot \| \nabla f(x_k)\|^2
  \]

  Далее воспользуемся трюком с суммой
  \[- \frac{c_1(c_2-1)C^2}{L} \cdot \sum\limits_{k=0}^n \| \nabla f(x_k)\|^2 \le \sum\limits_{k=0}^n [f(x_k)-f(x_{k+1})] = f(x_0) - f(x_{n+1}) \le f(x_0) - f^{*}
  \]
  Оценим:
  \[
  \sum\limits_{k=0}^n \|\nabla f(x_k)\|^2 \ge (n+1) \min_{0\le k\le n} \|\nabla f(x_k)\|^2
  \]

  Итого
  \[- \frac{c_1(c_2-1)C^2}{L} \cdot (n+1) \min_{0\le k\le n} \|\nabla f(x_k)\|^2 \le f(x_0) -f^{*}
  \]
  Что равносильно
  \[
  \min_{0\le k\le n} \| \nabla f(x_k)\| \le \sqrt{\frac{L (f(x_0) - f^{*})}{(n+1)c_1(1-c_2)C^{2}}}
  \]
  Частичные суммы $\sum_{k=0}^n \|\nabla f(x_k)\|^2$ ограничены одной константой, значит ряд сходится и $\|\nabla f(x_k)\| \rightarrow 0$.

  \begin{remark}
    Примечание про оптимальность. Если посмотреть на идею поиска оптимального шага (как в \hyperref[note:FFEA1F1D-B70C-4B4A-9386-8344A2A3BBC4]{Оценка убывания за шаг градиентного спуска}), то нам нужно будет максимизировать константу перед $\|\nabla f(x)\|^2$, чтобы аппроксимация была наиболее точной.
    Тогда нам нужно, чтобы $C \rightarrow 1 \Rightarrow \cos \theta_k \rightarrow -1$, а так как $0 < c_1 < c_2 < 1$, то $c_1(1-c_2) \rightarrow \frac{1}{4}$.
    Тогда оценка будет стремиться к таковой, при $c_1 = c_2 = \frac{1}{2}$ \\
    Получим, что
    \[
    \min_{0\le k \le n} \|\nabla f(x_k)\| \le \sqrt{\frac{4L (f(x_0) - f^{*})}{n+1}}
    \]
  \end{remark}
\end{proof}
